% ==============================================================================
% 08_limitations.tex
% Phụ trách: Thành viên 2 (Thực nghiệm, kết quả và kết luận)
% ==============================================================================

\section{Hạn chế của Nghiên cứu}
\label{sec:limitations}

Mặc dù quy trình thực nghiệm được thiết kế và kiểm toán nghiêm ngặt theo các tiêu chuẩn đóng băng dữ liệu (benchmark freeze), công trình này vẫn tồn tại một số giới hạn phương pháp luận cần được nhìn nhận một cách khách quan nhằm định hướng cho các nghiên cứu tiếp theo:

\paragraph{Phạm vi dữ liệu và ngôn ngữ (Corpus \& Language Scope).}
Toàn bộ thực nghiệm trong nghiên cứu này được tiến hành độc quyền trên kho tài liệu kỹ thuật của framework Angular bằng tiếng Anh, có cấu trúc cú pháp định dạng chuẩn Markdown. Do đó, các kết luận thực nghiệm không tự động suy diễn (generalize) cho các thể loại tài liệu phi cấu trúc khác (như văn bản tự do, tệp PDF đã qua xử lý OCR, báo cáo tài chính hay slide trình chiếu), cũng như các hệ thống tài liệu kỹ thuật đa ngôn ngữ. Cấu trúc phân cấp tiêu đề (heading hierarchy) và các khối cú pháp code block của Markdown có thể tạo ra những tương tác đặc thù với bộ phân tích cú pháp AST mà các định dạng tài liệu khác không sở hữu.

\paragraph{Độ bao phủ của Benchmark (Benchmark Coverage).}
Mặc dù kho ngữ liệu gốc bao gồm 384 tài liệu Markdown hoàn chỉnh, bộ benchmark đóng băng 140 cặp câu hỏi QA và 233 đơn vị gold evidence hiện tại mới chỉ khai thác và neo giữ bằng chứng trên 55 tài liệu nguồn được chọn lọc. Bên cạnh đó, dù 140 câu hỏi được phân bổ cân bằng trên 3 phân tầng độ khó chính (60 Easy, 40 Medium, 40 Hard), một số phân nhóm câu hỏi nhỏ theo dạng thức kỹ thuật sâu (question-type strata như cú pháp \texttt{syntax} với $n=1$, cơ chế bảo mật \texttt{security\_mechanism} với $n=1$, hay nguyên nhân--hệ quả \texttt{cause\_effect} với $n=2$) có kích thước mẫu quá nhỏ. Các kết quả trên những nhóm này thuần túy mang tính mô tả khám phá và chưa đủ lực thống kê để khái quát hóa thành kết luận độc lập.

\paragraph{Cấu hình pipeline và ngân sách đơn lẻ (Single-Configuration Pipeline).}
Nghiên cứu áp dụng thiết kế thực nghiệm kiểm soát nhằm cô lập tác động thuần túy của chunking strategy, do đó đã cố định một cấu hình duy nhất: một mô hình dense embedding (\texttt{text-embedding-3-small}), một mô hình sinh câu trả lời (\texttt{gpt-5-mini}), một độ sâu danh sách ứng viên ($k=50$), một ngưỡng retrieval token budget ($2.048$ tokens) và một mức context budget ($4.096$ tokens). Nghiên cứu chưa thực hiện kiểm thử lưới đa ngân sách (multi-budget grid search, ví dụ $1.024$ hay $4.096$ tokens), chưa tích hợp tầng tái xếp hạng học máy (cross-encoder reranker), và chưa khảo sát sự tương tác của các chiến lược phân đoạn với các kiến trúc mô hình ngôn ngữ lớn khác (như Claude, Gemini hay các mô hình open-weight như Llama).

\paragraph{Phương pháp đánh giá câu trả lời (Answer Evaluation Scope).}
Chất lượng câu trả lời sinh ra hiện được đo lường thông qua các chỉ số trùng khớp từ vựng cấp độ token (Token Precision, Token Recall, Token F1) đối chiếu với câu trả lời chuẩn tắc của chuyên gia sau khi chuẩn hóa Unicode NFKC. Dù phương pháp này đảm bảo tính tất định và khả năng tái lập độc lập 100\%, bản chất từ vựng của Token F1 có thể đánh giá chưa thỏa đáng các trường hợp mô hình diễn giải lại (paraphrase) câu trả lời bằng văn phong tự nhiên mà vẫn đảm bảo tính chính xác tuyệt đối về mặt kỹ thuật. Nghiên cứu chưa mở rộng sang việc áp dụng các khung đánh giá ngữ nghĩa đa chiều như LLM-as-a-judge có đối chứng hoặc human evaluation trên quy mô lớn.

\paragraph{Giới hạn của Prompt-Based Chunking Artifact.}
Chiến lược phân đoạn dựa trên prompt trong nghiên cứu này sử dụng dữ liệu được trích xuất từ artifact lịch sử cố định (\texttt{prompt\_based\_v1}) theo cam kết của phiên bản benchmark đóng băng v2. Nghiên cứu chưa khảo sát việc tối ưu hóa prompt động (dynamic prompt tuning), cung cấp ví dụ mẫu theo ngữ cảnh (few-shot in-context demonstrations) hay huấn luyện mô hình phân đoạn chuyên biệt (dedicated chunker fine-tuning). Do đó, kết quả của Prompt-based trong báo cáo phản ánh hiệu năng của artifact cụ thể này chứ không đại diện cho giới hạn tối đa của hướng tiếp cận prompt-based nói chung.
